\section{Diseño de la aplicación}

Esta sección describe el diseño general de la aplicación QFotoPaint6.

\subsection{Arquitectura}

- \textbf{Interfaz gráfica:} implementada con Qt (widgets y acciones en \texttt{mainwindow.*}).
- \textbf{Procesamiento de imagen:} funciones agrupadas en el módulo \texttt{imagenes.h/cpp}, usando OpenCV para operaciones sobre \texttt{cv::Mat}.
- \textbf{Vídeo y captura:} funcionalidades específicas en \texttt{video.h/cpp} que usan OpenCV (\texttt{VideoCapture}, \texttt{VideoWriter}).
- \textbf{Comunicación Qt / OpenCV:} conversiones entre \texttt{cv::Mat} y \texttt{QImage} para integrar con la API de Qt (por ejemplo para el portapapeles y visualización).

\subsection{Organización del código}

- Ficheros principales de UI: \texttt{mainwindow.cpp/.h} y el recurso \texttt{recursos.qrc}.
- Funcionalidades de imagen: \texttt{imagenes.h} (interfaz) y \texttt{imagenes.cpp} (implementación).
- Funcionalidades de vídeo/cámara: \texttt{video.h} y \texttt{video.cpp}.
- Ventanas de diálogo y utilidades: ficheros como \texttt{dialognueva.*}, \texttt{capturarvideo.*}, etc.

\subsection{Decisiones de diseño relevantes}

- Se mantiene un array estático \texttt{foto[MAX\_VENTANAS]} para gestionar las ventanas/imagenes abiertas, simplificando el control y la numeración de ventanas.
- Para integrar con Qt y el portapapeles se realizan conversiones controladas a \texttt{QImage} copiando los datos cuando es necesario para evitar problemas de propiedad de memoria.
- Se emplean callbacks de OpenCV (HighGUI) para manejar eventos de ratón sobre las ventanas de imagen, lo que mantiene la lógica de interacción separada de la GUI de Qt.
